\section{Introdução}
\label{sec:introduction}

Catálogos de aproximação costumam ser descritos por dois números: quantos
termos são usados e quão pequeno é o erro. Essa descrição omite um terceiro
recurso que se torna decisivo quando a representação precisa ser armazenada,
transmitida ou selecionada exatamente: o custo aritmético dos parâmetros. Uma
reta com coeficientes reais arbitrários pode ter descrição finita apenas em um
modelo idealizado; uma reta racional é finitamente codificável, mas o tamanho
de numeradores e denominadores pode crescer independentemente do número de
peças. Este trabalho torna essa complexidade aritmética uma variável de
primeira classe. Construímos uma biblioteca enumerável de segmentos afins
racionais, ancorados em células diádicas, e estudamos simultaneamente
endereçamento, aproximação, altura, comprimento binário e seleção exata.

O ponto de partida é propositalmente concreto. Cada átomo possui um centro
diádico, um valor racional e uma inclinação racional não nula. O endereço
aritmético descompacta esses dados, enquanto a interface global
$(a,b,[c,d])$ permite recuperar a célula de origem. A condição $b\ne0$ é
essencial para essa recuperação: uma peça constante colapsaria o intervalo de
imagem e apagaria o suporte. Ao impor continuidade nos nós compartilhados,
obtêm-se montagens lineares por partes que podem ser convertidas exatamente em
redes ReLU rasas. Essa estrutura fornece um laboratório no qual representação,
codificação e aproximação são verificáveis em aritmética racional de ponta a
ponta.

\subsection{Vizinhança bibliográfica e escopo}

A aproximação linear por partes pertence à teoria clássica de splines, de nós
livres e de aproximação não linear
\citep{deboor1973variable,burchard1974,deboor2001,schumaker2007,devore1998}.
A constante local $1/8$ da interpolação afim, a densidade em $C(K)$ e as taxas
de ordem dois para funções $C^2$ não são reivindicadas como novas. O mesmo vale
para a representação de splines lineares por expansões em hinges. Na teoria de
redes, resultados gerais relacionam profundidade, conectividade, pesos
codificáveis e classes de suavidade
\citep{petersenvoigtlaender2018,elbrachter2021,bolcskei2019,guhringraslan2021,
devorehaninpetrova2021}. Nosso incremento não é uma nova taxa abstrata de
aproximação neural; é a contabilidade explícita que resulta de combinar uma
malha diádica recuperável com coeficientes racionais de altura limitada.

Essa distinção muda perguntas aparentemente equivalentes. A cardinalidade de
um catálogo truncado não determina o maior índice produzido por uma enumeração
específica; um fluxo decodificável não tem necessariamente o mesmo comprimento
que a lista ingênua dos coeficientes; equivalência funcional entre uma montagem
e uma rede não preserva automaticamente o tamanho binário de toda
parametrização. A literatura de entropia métrica e codificação em árvores
oferece o parâmetro de comparação apropriado
\citep{kolmogorov1959,cohen2001,elias1975,coverthomas2006}. Aqui esses recursos
são separados por definição e depois ligados por teoremas com decodificadores
explícitos.

Há ainda duas vizinhanças distintas. A primeira é algorítmica: o princípio de
otimalidade de Bellman e a segmentação por programação dinâmica são clássicos
\citep{bellman1957,bellman1961,stone1961,perezvidal1994}. O papel das
recorrências deste trabalho é estabelecer uma teoria de exatidão para o
catálogo aritmético: a truncagem de altura torna finito o alfabeto de valores e
os dois valores de fronteira formam um estado suficiente para continuidade.
As pequenas instâncias executáveis são auditorias dessas identidades, não um
estudo de escalabilidade ou uma alegação de superioridade empírica. A segunda
vizinhança é funcional-analítica: classes de aproximação e imersões para redes
foram desenvolvidas em \citet{gribonval2022,grohs2023}. Quando se acrescenta a
autorregularidade do aproximante ao erro, surge um funcional de realização no
sentido de \citet{hristovivanov1990}; sua desigualdade inversa é universal. O
conteúdo específico da biblioteca está no lado direto com altura racional
controlada, e não na identificação abstrata do espaço de Besov.

Também delimitamos expressamente o que não é afirmado. A enumeração por altura
é minimax apenas para o calibre definido neste artigo; outro acoplamento entre
resolução, translação e coeficientes muda o expoente de contagem. As cotas de
programação dinâmica são pseudo-polinomiais nos orçamentos numéricos e as
experiências distribuídas são testes de regressão. O funcional de taxa bruto
não é uma quase-norma sobre escalares reais, pois o catálogo racional não é um
cone. Por fim, a caracterização pelo funcional de realização mede suavidade do
alvo para qualquer família suficientemente rica; não deve ser apresentada como
uma propriedade exclusiva da biblioteca.

\subsection{Contribuições demonstradas}

As contribuições próprias podem ser enunciadas sem depender do volume do
documento.

\begin{enumerate}[label=(\roman*),leftmargin=*]
 \item \textbf{Endereçamento recuperável.} Uma enumeração sem repetições
 produz todos os jatos $(x_0,y_0,s)$ com centro diádico não inteiro, valor
 racional e inclinação racional não nula. A interface global recupera
 exatamente a célula e, por consequência, o único endereço do átomo
 (\cref{lem:cell-recovery,thm:injective}).

 \item \textbf{Contagem e ótimo minimax sob altura.} Para o calibre fixado,
 há $\Theta(H^6)$ átomos de altura no máximo $H$. O gerador aritmético
 $\Gamma$, cuja inversa tem forma fechada, atinge índices de ordem
 $\Theta(H^8)$; a enumeração por cascas $\Gamma^\sharp$ atinge exatamente o
 menor maior índice possível. O fator $4/3$ entre os comprimentos assintóticos
 das palavras é demonstrado, não inferido
 (\cref{prop:gamma-exact-order,thm:sharp-optimal}).

 \item \textbf{Semântica e codec separados.} $\Gamma$ permanece a semântica
 operacional do átomo, enquanto $\Gamma^\sharp$ funciona como codec minimax.
 A conversão em ambos os sentidos preserva exatamente a célula e a função
 afim selecionada (\cref{rem:which-address}).

 \item \textbf{Serializações executáveis.} Três formatos prefix-free
 distinguem a escrita atômica, a forma ReLU com diádicos estruturais e a
 expansão racional ingênua. As duas primeiras usam $O(N\log H)$ bits; existe
 uma família para a qual a terceira exige $\Omega(NH)$, embora realize a mesma
 função (\cref{thm:coding-formats}).

 \item \textbf{Orçamento simultâneo e certificação.} Em bolas limitadas de
 $C^2[0,1]$, uma construção canônica alcança erro $O(N^{-2})$, uma agenda
 suficiente de altura $O(N^2)$ e descrição $O(N\log N)$. Para funções dadas por uma gramática
 elementar, o refinamento adaptativo usa aritmética intervalar racional com
 falha segura; assim, o majorante local é um certificado, não uma estimativa
 amostral (\cref{thm:joint-budget,thm:adaptive}).

 \item \textbf{Constantes estritamente diádicas.} Para interpolação em
 árvores diádicas, provamos o envelope assintótico
 $[I(f)^2/8,I(f)^2/2]$. Curvatura constante produz a lei exata por fases e o
 conjunto completo de acumulação $[A,4A]$; quadráticas por partes exibem uma
 função de fase periódica não constante
 (\cref{thm:general-dyadic-envelope,thm:quadratic-dyadic-constant,thm:piecewise-quadratic-count}).

 \item \textbf{Altura e regularidade.} Os pisos aritmético
 $H(N)\gtrsim N^s$ e geométrico $H(N)\gtrsim N$ são separados. O teorema de
 Jackson racional é fortalecido para uma versão uniforme sobre bolas de Besov,
 com constantes independentes da função e sem limiar individual
 (\cref{lem:rational-jackson-height,cor:uniform-besov-ball-embedding}).

 \item \textbf{Recíproca aguda na biblioteca original.} Para erro bruto
 $N^{-\alpha}$, crescimento Lipschitz $N^\gamma$ e $0<\alpha\le\gamma$, a ordem inversa
 $\alpha/(\alpha+\gamma)$ é ótima. A testemunha por pente de tendas é
 transferida da extensão marcada para a biblioteca original por uma
 perturbação diádica de inclinações não nulas. Para o funcional de realização,
 o lado direto preserva a agenda de altura e a direção inversa é identificada
 explicitamente como universal
 (\cref{thm:critical-tent-comb,rem:realization-functional-scope}).
\end{enumerate}

\subsection{Organização do documento-sistema}

As \cref{sec:indexing,sec:generator} constroem a enumeração e a recuperação da
célula. As \cref{sec:assemblies,sec:rates,sec:dyadic-trees,sec:adaptive}
tratam montagens, aproximação e refinamento certificado. A
\cref{sec:complexity} separa endereço, contagem, serialização, orçamento e
entropia. O bloco das \cref{sec:interface,sec:numerical} é preservado como
teoria auxiliar de seleção exata. As seções seguintes determinam as constantes
diádicas, introduzem os estratos de altura e estabelecem os resultados diretos
e inversos. A \cref{sec:verification} documenta a verificação computacional.

Esta versão unificada é mantida como documento-sistema porque algumas conexões
atravessam os blocos: a obstrução de translação reaparece tanto na seleção
quanto nas agendas de altura, e a distinção entre semântica e codec liga o
gerador à serialização. Para publicação em extensão convencional, o material
se decompõe naturalmente em peças sobre (a) constantes diádicas, (b)
endereçamento e codificação e (c) estratos de altura e teoremas inversos, com a
programação dinâmica como apêndice de exatidão. Essa arquitetura não remove
nenhum resultado do documento unificado.

Contagens estruturais da programação dinâmica usam inicialmente aritmética
racional de custo unitário; a \cref{thm:bit-complexity} fecha separadamente o
crescimento dos operandos e a complexidade de bits. Comprimento de
representação, custo de seleção e custo de decodificação permanecem recursos
distintos em todos os enunciados.
